\section{Introduction}
\label{sec:intro}

The rapid improvement of diffusion-based image generators~\cite{ho2020ddpm,rombach2022ldm} has made pixel-space authenticity cues unreliable: for many generators, no human --- and, increasingly, no naive classifier --- can distinguish a generated photograph from a real one on the basis of semantic content alone. This has reopened an old question in image forensics: if one looks not at what the image depicts but at the \emph{statistical structure} of its spatial frequencies, can generative artifacts still be recovered?

For GAN-generated imagery, the answer was famously yes. Wang \etal~\cite{wang2020cnn} showed that a single CNN trained on images from one GAN generalises surprisingly well to many others, and Durall \etal~\cite{durall2020watch} attributed this generalisation to consistent spectral distortions introduced by the transposed-convolution upsampling used in GAN decoders. Azimuthally averaged power spectra of real photographs approximately obey a $1/f^{\alpha}$ falloff~\cite{vanderschaaf1996}, and GANs systematically violate this law in detectable ways.

Diffusion models are architecturally different. They denoise iteratively in a latent space and decode through a VAE; their upsampling paths, noise schedules, and attention mechanisms are not the same as those of GANs. Whether the \emph{same frequency-domain story} survives the transition from GANs to diffusion models is an empirical question --- one that recent work has begun to address~\cite{corvi2023intriguing,ricker2024aeroblade,spai2025}, but whose answer is not yet settled, especially when evaluation is confounded by JPEG-induced artifacts and dataset biases~\cite{fakeorjpeg2024,synthbuster2023}.

This paper takes that question seriously and asks a sharper version of it: \emph{does the diffusion-era fingerprint live in the isotropic component of the spectrum, or does it require the full 2D structure?} The azimuthal average is a rotation-invariant summary that throws away all directional information; the full 2D spectrum retains grid patterns, axis-aligned biases, and other anisotropic artifacts. Prior frequency-domain detectors have chosen one representation or the other, often for engineering reasons; to our knowledge, no paper has run a controlled head-to-head comparison grounded in this isotropy hypothesis.

\paragraph{Contributions.} We make three contributions.
\begin{enumerate}[itemsep=1pt,topsep=1pt]
    \item \textbf{A controlled 1D--vs--2D frequency-only comparison.} We train two architecturally matched CNNs on the same images, differing only in whether they see the full 2D log-power spectrum or its azimuthally averaged 1D profile. The 1D model has $\sim$180k parameters and the 2D model $\sim$4M; neither uses pretrained weights, semantic features, or data augmentation. The comparison directly tests whether diffusion artifacts are isotropic (in which case the 1D model should suffice) or anisotropic (in which case the 2D model should dominate as generator diversity grows).
    \item \textbf{Empirical evidence on two diffusion-era benchmarks.} On CIFAKE~\cite{cifake} (Stable Diffusion v1.4 vs CIFAR-10) our 2D detector reaches 0.953 validation AUC and our 1D detector 0.940; on Defactify~\cite{defactify} (a multi-generator mixture) the numbers are 0.905 and 0.844. The gap widens from 0.013 AUC on a single-generator benchmark to 0.061 AUC on a multi-generator one, consistent with anisotropic artifacts becoming more important as the generator set diversifies.
    \item \textbf{A planned cross-generator and adversarial robustness study.} We describe, and motivate, a still-running set of experiments on the full eight-generator GenImage~\cite{genimage} training benchmark and on the GenImage++~\cite{genimagepp} holdout, together with a Fourier-space PGD attack~\cite{madry2018pgd} that perturbs the input in frequency coordinates directly. The attack is intended both as a robustness test and as an interpretability tool: the frequencies modified by successful attacks should identify the spectral bands the detector is actually keying on.
\end{enumerate}

We emphasise that the GenImage and adversarial experiments are \emph{in progress} at the time of submission; we include them here with their full experimental design and the outcomes we expect, so that the final paper can be read as a complete scientific argument rather than a retroactive reconstruction. Sections describing unfinished experiments are clearly marked.

\paragraph{Scope and constraints.} Throughout this work we restrict ourselves to features derived purely from the Fourier transform of single images. We use no pretrained visual encoders, no CLIP or DINO features, no metadata, and no semantic labels. This is a deliberate constraint: frequency-only detection is both cheap at inference time and more tractable to reason about than a black-box foundation-model detector, and we want to know how far it can be pushed on its own terms.
